% G. Discusión y conclusiones, incluyendo responsabilidad legal, ética y profesional
% (Guía apdo. 3.G). Extensión orientativa: 3-4 páginas.
\chapter{Discusión y conclusiones}
\label{cap:conclusiones}

\section{Discusión}

Los resultados obtenidos a lo largo del desarrollo y de las distintas campañas experimentales
permiten extraer varias reflexiones de fondo en torno al despliegue de sistemas de mantenimiento
predictivo basados en el paradigma de computación en el extremo (\textit{Edge Computing}).

En primer lugar, la viabilidad computacional de la inferencia en microcontroladores de bajo coste
ha quedado demostrada empíricamente. Frente a la tendencia habitual en la literatura reciente
de trasladar redes neuronales profundas mediante entornos de cuantización como TensorFlow Lite
Micro, el presente trabajo evidencia que modelos no paramétricos de detección de anomalías de
una clase (en particular, \textit{Local Outlier Factor}) combinados con una extracción rigurosa
de características en el dominio del tiempo y de la frecuencia ofrecen una solución altamente
competente y eficiente. Con un tiempo de ejecución mediano de \mbox{14 \textmu s} y un peor caso
inferior a \mbox{2 ms}, la inferencia consume únicamente el \mbox{0,00005~\%} del ciclo de trabajo
de \mbox{30 s} del nodo, demostrando que la restricción computacional en el borde queda superada
sin necesidad de recurrir a procesadores dedicados de señal o coprocesadores de inteligencia
artificial.

En segundo lugar, el análisis ha puesto de manifiesto que la integridad de la cadena de
adquisición condiciona por completo la validez del diagnóstico. Tal como se documentó durante la
fase de diseño y captura, las perturbaciones en el bus I2C derivadas de la vibración mecánica
son capaces de introducir muestras corruptas aisladas que disparan la kurtosis temporal y
generan picos espectrales espurios. Estos artefactos numéricos reproducen de forma artificial la
firma espectral de un fallo sobre una máquina completamente sana. De ello se deriva un principio
metodológico fundamental: los mecanismos de verificación de integridad (comprobación de
continuidad, recuento de reintentos de bus y filtrado de calidad) no constituyen un mero
preproceso descartable, sino que forman parte indisoluble del algoritmo de diagnóstico. La
decisión del nodo de emitir un veredicto de «no evaluable» ante ráfagas degradadas evita
decisiones erróneas basadas en datos corruptos.

En tercer lugar, la integración multimodal ha demostrado ser un elemento diferenciador. La
presencia del canal acústico permitió corroborar de manera independiente y sin compartir cadena
de sensado la familia armónica detectada en el activo con fallo (campaña EXP-003), donde el
reparto energético en las bandas de frecuencia se invirtió respecto a la condición nominal.
Asimismo, el uso de variables térmicas diferenciales (\texttt{dif\_termico}) aporta robustez
frente a variaciones lentas de la temperatura ambiente exterior.

Por último, la introducción de una etapa de histéresis temporal basada en rachas consecutivas de
anomalías (\texttt{streak = 3}) ha demostrado ser indispensable para la viabilidad operativa del
sistema. En entornos industriales y domésticos, los transitorios de arranque y las pequeñas
perturbaciones mecánicas aisladas generan ráfagas anómalas esporádicas. La exigencia de tres
veredictos anómalos consecutivos reduce en dos órdenes de magnitud el número de avisos frente a
la fracción de ráfagas marcadas.

Procede sin embargo acotar su eficacia con lo observado. La histéresis absorbe las desviaciones
aisladas, pero \textbf{no distingue el origen de una desviación sostenida}: en la campaña
EXP-008 emitió un aviso ante un transitorio acústico de noventa segundos provocado por la
apertura de una puerta, que es un suceso operativo normal y no un fallo. Sobre la captura
completa se emitieron dieciséis avisos en \mbox{22,80 h} de operación con el activo en estado
nominal, cifra que constituye el límite práctico del sistema en su configuración actual y no un
resultado satisfactorio.

\section{Conclusiones}

El presente Trabajo de Fin de Máster ha abordado el diseño, desarrollo, despliegue y validación
experimental de un sistema modular IoT para el diagnóstico predictivo de maquinaria
electromecánica en el borde de la red. A continuación se evalúa el grado de consecución de cada
uno de los objetivos específicos planteados al inicio del proyecto:

\begin{enumerate}
    \item \textbf{Definición del escenario de prueba:} \textit{Alcanzado.} Se seleccionó como
        activo representativo un compresor hermético de refrigeración doméstica, caracterizando
        su régimen nominal de giro a \mbox{49,15 Hz} y estableciendo los límites temporales y
        espectrales de referencia a partir de más de 22 episodios de marcha.
    \item \textbf{Diseño de la arquitectura del sistema:} \textit{Alcanzado.} Se definió una
        arquitectura Edge--Fog desacoplada en tres canales MQTT independientes (canal lento de
        sensores, canal de ráfagas espectrales y canal de estado del nodo). Esta estructura
        permite reducir el tráfico de red en un orden de magnitud al transmitir únicamente el
        veredicto de salud (\texttt{fridge/status}, \mbox{124 B}) hacia los sistemas de nivel
        superior.
    \item \textbf{Selección e integración de hardware y software:} \textit{Alcanzado.} Se integró
        con éxito un microcontrolador ESP32-S3 junto con un acelerómetro triaxial MPU-6050, una
        sonda térmica DS18B20 y un micrófono digital INMP441. Se desarrollaron rutinas de
        recuperación automática de bus y extracción de características en C++ (FFT de 1024 puntos
        a \mbox{1 kHz}) sin dependencias externas.
    \item \textbf{Desarrollo del motor de análisis inteligente:} \textit{Alcanzado.} Se formuló
        un modelo de detección de anomalías de una clase basado en \textit{Local Outlier Factor}
        (LOF), ajustado mediante un protocolo de partición cronológica por episodios libre de
        sesgo de espionaje. Alcanza una tasa de detección del 100~\% sobre el único fallo real
        observado y un 7,8~\% de falsos positivos sobre episodios no vistos. Un análisis de
        sensibilidad frente a cinco direcciones de fallo indica que el modelo reaccionaría a
        todas ellas, si bien \textbf{ese análisis emplea perturbaciones sintéticas del conjunto
        nominal y no constituye evidencia experimental}.
    \item \textbf{Validación experimental:} \textit{Parcialmente alcanzado.} El sistema se
        contrastó frente a un fallo armónico real en un activo secundario (EXP-003), con verdad
        de referencia externa a la instrumentación, y el detector lo identificó en la totalidad
        de las ráfagas. \textbf{No se ha inducido ningún fallo sobre el activo de referencia}, de
        modo que la limitación de mayor peso del trabajo permanece: el contraste se establece
        entre máquinas distintas y no separa el efecto del fallo del de la variabilidad entre
        ejemplares.

        La campaña EXP-008, consistente en abrir la puerta del aparato, no suple esa carencia.
        El activo se encontraba sano y una prueba de ablación establece que el aviso lo provocó
        el desplazamiento de la energía acústica y no la firma vibratoria ni la térmica: anulando
        el canal acústico el episodio deja de detectarse, y anulando cualquier otro grupo de
        características sigue detectándose. Abrir una puerta no degrada el compresor, de modo
        que el episodio constituye un falso positivo sobre un suceso operativo normal.

        Sí aportó, en cambio, dos resultados de valor: la constatación de que el canal acústico
        es el responsable de \mbox{8,3} puntos de la tasa de ráfagas marcadas, y la observación
        de que el efecto térmico sostenido de la perturbación —un incremento de \mbox{5,5 °C} en
        el diferencial— fue señalado por el detector con treinta minutos de retraso respecto del
        transitorio acústico.
    \item \textbf{Pruebas en entorno real:} \textit{Alcanzado.} El nodo completo con inferencia
        embarcada fue desplegado en operación continua durante \mbox{22,80 h} consecutivas
        (EXP-007), acumulando 2659 veredictos con una paridad del 99,42~\% frente al recálculo
        en el equipo de análisis —1030 de 1036 ráfagas evaluables— y una mediana de inferencia
        de \mbox{14 \textmu s}. Debe consignarse que la fracción de ráfagas marcadas como
        anómalas sobre el activo sano fue del \mbox{20,4 \%}, muy por encima del \mbox{7,8 \%}
        estimado en prueba, y que esa discrepancia \textbf{no está explicada}: ninguna
        característica la justifica por sí sola.
\end{enumerate}

Como balance global, el trabajo establece de forma cuantitativa que es factible desplegar
algoritmos de aprendizaje automático no supervisado en microcontroladores de bajo coste: la
inferencia ocupa el \mbox{0,004 \%} del ciclo y reproduce el veredicto de la implementación de
referencia. Lo que \textbf{no} establece, y conviene enunciarlo con la misma claridad, es que el
detector resultante sea utilizable en explotación: dieciséis avisos en \mbox{22,80 h} sobre un
activo sano lo desaconsejan, y la especificidad del sistema queda pendiente de la reducción de
esa tasa.

\section{Consideraciones legales, éticas y profesionales}

El sistema desarrollado incorpora un micrófono digital, lo que exige delimitar estrictamente su
ámbito de aplicación y tratamiento de datos. La señal acústica se procesa localmente en el nodo
en bloques de corta duración para obtener exclusivamente un vector agregado de energía por bandas
espectrales, descartándose de inmediato la señal de audio en bruto. El sistema no almacena ni
transmite grabaciones sonoras, garantizando la imposibilidad técnica de reconstruir
conversaciones o identificar personas. Esta decisión de diseño garantiza el cumplimiento del
principio de minimización de datos establecido en el Reglamento (UE) 2016/679 (RGPD)~\cite{gdpr}.

En el ámbito de la ciberseguridad, la prueba de concepto evaluada opera en una red Wi-Fi local
dedicada generada por el concentrador Raspberry Pi, con aislamiento respecto a redes públicas.
No obstante, el traslado de este diseño a un entorno industrial de producción requiere
incorporar mecanismos de autenticación mutua mediante certificados X.509 y cifrado TLS en el
broker MQTT, evitando la interceptación de telemetría y la inyección de veredictos maliciosos.

Desde el punto de vista profesional y de responsabilidad de ingeniería, debe subrayarse que un
sistema de monitorización IoT de bajo coste está concebido como una herramienta de alerta
temprana y apoyo al mantenimiento, y en ningún caso exime del cumplimiento de los planes de
mantenimiento preventivo e inspección reglamentaria exigibles por la normativa técnica de
seguridad industrial.

\section{Líneas futuras}

A partir de los resultados alcanzados y de las limitaciones identificadas, se proponen las
siguientes líneas de trabajo futuro:

\begin{itemize}
    \item \textbf{Fallo inducido sobre el activo de referencia.} Es la primera línea por orden
        de importancia, dado que resuelve la limitación de mayor peso del trabajo. Un
        aflojamiento del anclaje o un desequilibrio de masa alteran la firma vibratoria del
        propio activo y permiten separar el efecto del fallo del de la variabilidad entre
        ejemplares, cosa que el contraste entre máquinas distintas no permite.
    \item \textbf{Reducción de la tasa de avisos falsos.} Los datos sugieren una vía que no se
        ha ensayado: exigir que la desviación se manifieste también en el canal de vibración
        antes de notificar, con lo que el canal acústico conservaría su función confirmatoria
        —fue el que aportó la evidencia independiente del fallo real— sin poder disparar por sí
        solo ante una perturbación del entorno.
    \item \textbf{Aprendizaje autónomo en el borde (\textit{On-Device Learning}):} Desarrollar
        rutinas de autocalibración en el firmware para que el nodo estime de forma autónoma sus
        medianas nominales durante las primeras horas de puesta en marcha, eliminando el ciclo
        de exportación de cabeceras desde el servidor.
    \item \textbf{Ensayos de degradación térmica controlada:} Ejecutar campañas prolongadas de
        obstrucción parcial de condensador y sobrecarga de puerta abierta para evaluar la
        dinámica de respuesta del gradiente térmico a lo largo de varios ciclos día--noche.
    \item \textbf{Generalización sobre flotas de activos:} Desplegar nodos en múltiples unidades
        de compresores heterogéneos para analizar la transferibilidad del modelo y la
        variabilidad de los umbrales de vibración entre distintos fabricantes.
    \item \textbf{Ciberseguridad robusta:} Implementar soporte TLS/MQTTS nativo en el ESP32-S3
        aprovechando los aceleradores criptográficos por hardware del chip.
    \item \textbf{Ampliación de la resolución espectral:} Aumentar el número de picos armónicos
        publicados por el firmware para permitir el seguimiento simultáneo de órdenes
        superiores de giro (por encima del tercer pico).
\end{itemize}
